\section{La forme normale de Jordan}
\newdef{Un \emph{bloc de Jordan} est définie comme
		\[	
		J(\lambda,n)=\begin{pmatrix}
			\lambda & 1 &  &  &  \\
			& \lambda & 1 &  & \\
			&  &  \ddots &  \ddots\\
			&  &  &  \lambda & 1 \\
			&  &  &  & \lambda
		\end{pmatrix} \in \C^{n \times n}
	\]
	où $\lambda \in \C$ et $n \in \N$. On remarque que $J(\lambda, n) = \mathrm{diag}(\lambda, \ldots, \lambda) + J(0,n)$ où $J(0,n)$ est une matrice nilpotente, puisque $J(0, n)^{n} = 0$. De plus on dit qu'une matrice $A \in \C^{n \times n}$ est une \emph{forme normale de Jordan} si
	\[
	A = \begin{pmatrix} J(\lambda_1, n_1) & & \\ & \ddots &  \\ &  & J(\lambda_k, n_k) \end{pmatrix}
	\]
	avec $\sum n_k = n$.
}{}

\newlem{Soit $A \in \C^{n \times n}$ une matrice normale de Jordan alors
	\[
	\chi_A(x) = (-1)^n \prod_{i=1}^k (x - \lambda_i)^{n_i}
	\]
}{}

\prf{Comme $A - x \mathrm{Id}$ est une matrice triangulaire supérieure, alors son déterminant est donné par le produit des termes diagonaux, ainsi on a bien
	\[
	\chi_A(x) = \det(A - x \mathrm{Id}) = (-1)^n \prod_{i=1}^n (x - \lambda_i)^{n_i}
	\]
}



\newlem{Soit $J \in \C^{n \times n}$ une forme normale de Jordan. Alors
	\[
	 J=\underbrace{ \begin{pmatrix}
	 		\lambda_{1}I_{n_{1}} &  &  \\
	 		& \ddots &  \\
	 		&  & \lambda _{k}I_{n_{k}}
	 \end{pmatrix} }_{ D }+\underbrace{ \begin{pmatrix}
	 		J(0,n_{1}) &  &  \\
	 		& \ddots &  \\
	 		&  & J(0,n_{k})
	 \end{pmatrix} }_{ N }.
	\]
	où $D$ et $N$ commutent, et $N$ est nilpotent.
}{}

\prf{On a
	\[
	N^i = \begin{pmatrix}
		J(0,n_{1})^i &  &  \\
		& \ddots &  \\
		&  & J(0,n_{k})^i
	\end{pmatrix}
	\]
	Soit $j = \max\{n_1, \ldots, n_k\}$, alors $N^j = 0$. On a que $D$ et $N$ commutent
	\[
	\lambda_i I_{n_i} J(0, n_i) = J(0, n_i) \lambda_i I_{n_i} \implies D N = ND
	\]
}
\subsection{Déterminer la forme normale de Jordan}
\newthm{Soit $A \in \C^{n \times n}$, alors il existe une matrice $P \in \C^{n \times n}$ inversible et une forme normale de Jordan $J \in \C^{n \times n}$ tel que
	\[
	A = P J P^{-1}
	\]
}{Réduction de Jordan}

\newlem{Soit $J \in \C^{n \times n}$ une forme normale de Jordan et soient $\lambda_1, \ldots, \lambda_k \in \C$ des valeurs propres distincts de $A \in \C^{n \times n}$. Alors le polynôme minimal de $J$ est 
	\[
	\mu_j(x) = (x- \lambda_1)^{e_i} \cdot \ldots \cdot (x- \lambda_k)^{e_i}
	\]
	où $e_i = \max \{n_i \mid J(\lambda_i, n_i) \text{ est un bloc de Jordan}\}$.
}{}

\prf{Soit 
	\[
	J=\begin{pmatrix}
		J(\lambda_{1},n_{1}) \\
		& J(\lambda_{1},n_{2}) \\
		&  & \ddots  \\
		&  &  & J(\lambda_{1},n_{k}) \\
		&  &  &  &  J(\lambda_{2},n_{k+1}) \\
		&  &  &  &  & \\
		&  &  &  &  & \ddots 
	\end{pmatrix}
	\]
	alors en posant $e_1 = \max\{n_i \mid J(\lambda_1, n_i) \text{ est un bloc de Jordan}\} = \max \{ n_i \mid i \in \llbracket 1, k \rrbracket \}$, alors
		\[
		\begin{pmatrix}
			J(0,n_{1})^{e_{1}} \\
			&  \ddots \\
			&  &  J(0,n_{k})^{e_{1}}
		\end{pmatrix}=0
		\]
	et pareil pour les autres blocs. 
}	
\newlem{Soit $k \in \N$, alors
	\[
	\dim \ker(A- \lambda I)^k = \dim \ker (J - \lambda I)^k
	\]
}{}

\prf{Pour $\lambda \in \C$ on voit que
	\[
	(A - \lambda I)^k = (P J P^{-1} - \lambda P P^{-1})^k = (P(J - \lambda I)P^{-1})^k = P(J - \lambda I)^k P^{-1}
	\]
	Ainsi un élément $x \in \C$ est contenu dans le noyau de $(A - \lambda)^k$, si et seulement si $P^{-1} x \in \ker (J - \lambda I)^k$, d'où la proposition.
}


\newthm{Soit $A \in \C^{n \times n}$ et $\lambda_1, \ldots, \lambda_k \in \C$ les valeurs propres distinctes de $A$. Soit
	\[
	d_i = \dim \ker(A - \lambda_1 I)^i
	\]
	alors
	\[
	n_i = d_i - d_{i-1}
	\]
	correspond au nombre de blocs de Jordan dont $\lambda_1$ est l'élément diagonal et de longueur aux moins $i$.
}{}

\prf{Soit $J$ une forme de Jordan de $A$, alors $J$ est de la forme
	\[
	J = \mathrm{diag} (J(\lambda_1, n_{1, 1}), J(\lambda_1, n_{2,1}), \ldots J(\lambda_1, n_{k,1}), J(\lambda_2, n_{2,1}), \ldots, J(\lambda_n, n_{n,m}))
	\]
	Ainsi par le lemme qui précede, on considère la quantité $(J - \lambda_1 I)^i$, ce qui nous donne
	\[
	(J - \lambda_1 I)^i = \mathrm{diag}  (J(0, n_{1, 1})^i, \ldots, J(0, n_{2,1})^i, \ldots J(0, n_{k,1})^i, J(\lambda_2 - \lambda_1, n_{2,1})^i, \ldots, J(\lambda_n - \lambda_1, n_{n,m})^i)
	\]
	On remarque alors que $J(0, n_{j,1})^i = 0$ si $i \geqslant n_{j,1}$ pour tout $j = 1,2, \ldots, k$. Donc on a bien
	\[
	d_i = d_{i-1} + \#\text{de bloc de longueurs} \geqslant i 
	\]
	d'où le théorème.
}
\newcor{Soit $A \in \C^{n \times n}$, $J \in \C^{n \times n}$ une forme normale de Jordan de $A$ et $\lambda \in \C$ une valeur propre de $A$. Soit $k \in \entiers$ et $n_k$ le nombre de bloc de Jordan avec $\lambda$ comme élément diagonal et de longueur $k$ est donné par
	\[
	n_k = 2 d_i - d_{i-1} - d_{i+1}
	\]
}{}

\subsection{Décomposition selon le polynôme caractéristique}
\newdef{Soient $W \subset \C^n$ un sous-espace vectoriel et soit $A \in \C^{n \times n}$, on dit que $W$ est \emph{invariant} sous $A$,
	\[
	\forall x \in W, \quad A x \in W
	\]
}{}

\newlem{Soient $f \in \K[X]$ et $A \in \K^{n \times n}$ une matrice quelconque, alors
	\[
	\ker (f(A))
	\]
	est invariant sous $A$.
}{}

\prf{Soit $x \in \ker(f(A))$ on trouve que
	\[
	f(A) \cdot Av = A \cdot f(A) \cdot x = A \cdot 0 = 0
	\]
	Ainsi $Av \in \ker(f(A))$.
}

\newthm{Soit $A \in \K^{n \times n}$ et soit $f(x) = f_1(x) \cdot f_2(x) \in \K[X]$ tel que
	\begin{enumerate}
		\item $\deg f_i > 0$ pour tout $i$ 	
		\item $\mathrm{pgcd}(f_1, f_2) = 1$
	\end{enumerate}
	Alors
	\[
	\ker(f(A)) = \ker(f_1(A)) \oplus \ker(f_2(A))
	\]
}{Lemme des noyaux}

\prf{Comme $\mathrm{pgcd}(f_1, f_2) = 1$, par Bézout on sait qu'il existe $g_1(x), g_2(x) \in \K[X]$ tel que
	\[
	1 = g_1(x) f_1(x) + g_2(x) f_2(x)
	\]
	et alors par évaluation
	\[
	I_n = g_1(A) \cdot f_1(A) + g_2(A) \cdot f_2(A)
	\]
	Ainsi pour $v \in \ker(f(A))$, on a
	\[
	v = g_1(A) \cdot f_1(A)v + g_2(A) \cdot f_2(A)v
	\]
	On remarque que
	\[
	g_1(A) \cdot f_1(A) v \in \ker (f_2(A))
	\]
	en effet
	\[
	f_2(A) g_1(A) f_1(A) v = g_1(A) f_1(A) f_2(A) v = g_1(A) f(A) v = 0
	\]
	il en va de même que $g_2(A) \cdot f_2(A) v \in \ker(f_1(A))$. Soit $v \in \ker(f_1(A)) \cap \ker(f_2(A))$, on a 
	\[
		v = g_1(A) \cdot f_1(A)v + g_2(A) \cdot f_2(A)v = 0
	\]
	ce qui démontre bien que la somme est directe.
}

\newcor{On peut sans problème généraliser ce résultat, ainsi soit $f(x) = f_1(x) \cdot \ldots \cdot f_k(x) \in \K[X]$ tel que
	\begin{enumerate}
		\item $\deg f_i > 0$ pour tout $i$
		\item $\mathrm{pgcd}(f_i, f_j) = 1$ pour tout $i \neq j$
	\end{enumerate}
	alors
	\[
	\ker(f(A)) = \ker(f_1(A)) \oplus \dots \oplus \ker(f_k(A)) = \bigoplus_{i=1}^k \ker(f_i(A))
	\]
}

\newthm{Soit $A \in \C^{n \times n}$, alors $\C^n$ est la somme directe de sous-espace vectoriel $\C^{n} = W_1 \oplus \dots \oplus W_k$ tels que
	\begin{enumerate}
		\item $W_i$ est invariant par $A$ pour tout $i$
		\item $A_{\mid W_i}$ est de la forme $\lambda_i I_{n_i} + N_i$ où $N_i$ est nilpotente
	\end{enumerate}
}{}

\prf{Soit $\chi_A(x)$ le polynôme caractéristique de $A$, par Cayley-Hamilton, on sait que $\chi_A(A) = 0$. De plus par le théorème fondamentale de l'algèbre on sait que $\chi_A$ se laisse factoriser sur $\C$ comme
	\[
	\chi_A(x) = (-1)^n (x - \lambda_1)^{n_1} \cdot \ldots \cdot (x - \lambda_k)^{n_k}
	\]
	avec des $\lambda_i$ deux à deux distincts. Alors par le théorème qui précède on a bien,
	\[
	\C^n = \ker(\chi_A(A)) = \ker(A - \lambda_1 I_n)^{n_1} \oplus \ldots \oplus \ker(A - \lambda_k I_n)^{n_k}
	\]
	On pose alors $W_i = \ker (A - \lambda_iI_n)^{n_i}$. De plus
	\[
	A_{\mid W_i} = (A - \lambda_i I_n)_{\mid W_i} + \lambda_i {I_n}_{\mid W_i} := N_i + \lambda_i I_n
	\]
	or $N_i$ est bien nilpotent car
	\[
	N_i^{n_i} = (A - \lambda_i I_n)_{\mid W_i}^{n_i} = 0
	\]
}







